\chapter{Conclusions and Further Work}

The thesis project establishes the foundations of a nonlinear flight dynamics simulation environment for a Boeing 777-like aircraft. The current work defines the coordinate conventions, state and input vectors, geometry data, mass and inertia estimates, standard atmosphere model, staged aerodynamic database, propulsion-load seed and validation-oriented reference cases. These elements provide the data layer and first force-and-moment adapters required before the full six-degree-of-freedom plant and NMPC controller can be integrated.

The most important modelling decision at the current stage is the separation of fixed aircraft properties from environment-dependent force quantities. Mass and inertia are stored as aircraft properties, while weight and thrust-to-weight quantities are computed using the current gravity assumption. Although the project presently uses fixed standard gravity, this interface leaves a clear path toward a future variable-gravity or higher-fidelity environment model.

Further work should focus on adding control-surface actuator dynamics, integrating the nonlinear equations of motion, extending trim routines to the full coupled plant and validating the open-loop dynamic behaviour. After these steps, the longitudinal and lateral-directional NMPC controllers can be implemented and tested in the planned simulation scenarios.

The final engineering objective is not to reproduce a certified Boeing 777 simulator, but to obtain a transparent, consistent and extensible platform for studying flight dynamics and constrained autopilot design.
